% Chunk: §7.7 Field Redefinitions and Redundant Couplings
% Source: Weinberg QFT Vol. I, printed pp. 331--332 (physical PDF pp. 354--355)
% Status: transcribed; source-reviewed; compile-clean; render-reviewed
% Semantic start: §7.7 heading on p. 354. End: paragraph before Appendix heading on p. 355.
% Figures: none. Uncertainties: none.

\subsection[Field Redefinitions and Redundant Couplings]{Field Redefinitions and Redundant
Couplings\texorpdfstring{\protect\footnotemark}{}}
\footnotetext{This section lies somewhat out of the book's main line of
development, and may be omitted in a first reading.}

Observables like masses and $S$-matrix elements are independent of some of the
coupling parameters in any action, known as the \emph{redundant parameters}.
This is because changes in these parameters can be undone by simply redefining
the field variables. A continuous redefinition of the fields, such as an
infinitesimal local transformation
$\Psi^\ell(x)\to\Psi^\ell(x)+\epsilon F^\ell(\Psi(x),
\partial_\mu\Psi(x),\ldots)$, clearly cannot affect any \emph{observable} of
the theory,\footnote{For instance, the theorem of Section 10.2 shows that as
long as we multiply by the correct field renormalization constants, $S$-matrix
elements can be obtained from the vacuum expectation value of a time-ordered
product of \emph{any} operators that have non-vanishing matrix elements
between the vacuum and the one-particle states of the particles participating
in the reaction.} though, of course, it would change the values of matrix
elements of the fields themselves.

How can we tell whether some variation in the parameters of a theory can be
cancelled by a field redefinition? A continuous local field redefinition will
produce a change in the action of the form
\begin{equation}
\delta I[\Psi]=\epsilon\sum_\ell\int d^4x\,
\frac{\delta I[\Psi]}{\delta\Psi^\ell(x)}
F^\ell(\Psi(x),\partial\Psi(x),\ldots).
\tag{7.7.1}\label{eq:7.7.1}
\end{equation}
So any change $\delta g_i$ in the coupling parameters $g_i$, for which the
change in the action is of the form
\begin{equation}
\sum_i\frac{\partial I}{\partial g_i}\delta g_i
=-\epsilon\sum_\ell\int d^4x\,
\frac{\delta I[\Psi]}{\delta\Psi^\ell(x)}
F^\ell(\Psi(x),\partial\Psi(x),\ldots),
\tag{7.7.2}\label{eq:7.7.2}
\end{equation}
may be compensated by a field redefinition
\[
\Psi^\ell(x)\to\Psi^\ell(x)+\epsilon F^\ell
(\Psi(x),\partial_\mu\Psi(x),\ldots),
\]
and therefore can have no effect on any observables. In other words,
\emph{a coupling parameter is redundant if the change in the action when we
vary this parameter vanishes when we use the field equations
$\delta I/\delta\Psi^\ell=0$}.

For example, suppose we write the Lagrangian density of a scalar field theory
in the form
\[
\mathcal L=-\frac12Z(\partial^\mu\Phi\partial_\mu\Phi+m^2\Phi^2)
-\frac1{24}gZ^2\Phi^4.
\]
The constant $Z$ is an redundant coupling, because
\[
\frac{\partial I}{\partial Z}=\frac12\int d^4x\,
\Phi\left(\Box\Phi-m^2\Phi-\frac16gZ\Phi^3\right),
\]
and this vanishes when we use the field equation
\[
\Box\Phi-m^2\Phi=\frac16gZ\Phi^3.
\]
On the other hand, neither the bare mass $m$ nor the coupling $g$ are
redundant, and no function of $m$ and $g$ is redundant.

In this example, the field redefinition needed to compensate for a change of
$Z$ is a simple rescaling, in which $F$ is proportional to $\Phi$. (For this
reason $Z$ is called a field renormalization constant.) This is the most
general field transformation that leaves the general form of this action
invariant. But for the more general actions considered in Sections 12.3 and
12.4, with arbitrary numbers of fields and derivatives, we would have to
consider non-linear as well as linear field redefinitions, and an infinite
subset of the parameters of the theory would be redundant.
